\chapter{Computational Results}
\label{chap:results}

\vspace*{\fill}

This chapter evaluates the hot-started nullspace active-set method on generated instances of
the \textbf{airline crew scheduling problem}. Every instance is solved in two variants: with
equality coverage rows only, and with 30\% of the coverage rows relaxed to ``$\ge$''
inequalities. The active-set method is run with both factorizations of the effective reduced
Hessian: the sparse factorization of \cref{chap:updates} and the dense factorization of
\cref{chap:dense}.

\medskip

We first check \textbf{correctness}: every column generation method must reproduce the
optimum of the full linear program, which is still small enough to be solved directly. We
then study the \textbf{convergence behaviour} of regularized column generation, which
consists of two clearly separated kinds of iterations, and the \textbf{computational time},
broken down by solver, by iteration and by operation inside the active-set method, up to
instances with almost a thousand linking rows.

\medskip

The last sections look inside the method. We follow the sizes of the nullspace sets over a run
with inequality rows, measure the \textbf{contents of the sparse factorization} and the value
of periodic refactorization, and finally compare the \textbf{two factorizations} operation by
operation, in memory and in accuracy.

\vspace*{\fill}

\newpage


\section{Test Problem: Airline Crew Scheduling}
\label{sec:test-problem}

\subsection{The model}

In the crew scheduling problem, $K$ pilots have to be assigned schedules that together
operate a set of $m$ flights (\cref{ex:crew}). The schedules of a pilot are the source--sink
paths in a \textbf{time--space network}: a directed acyclic graph $G = (V,E)$ whose nodes
represent states of a pilot at given times, and whose edges represent the allowed transitions
between them. The edge costs of pilot $k$ form the vector $c^{(k)}\in\R^{|E|}$, and every edge
operates a subset of the flights, recorded in the \textbf{task matrix}
$A_{\mathrm{task}}\in\{0,1\}^{m\times|E|}$. With $\xi_k\in\{0,1\}^{|E|}$ the edge vector of the
path of pilot $k$ and $F$ the node--edge incidence matrix of $G$ ($F_{\mathrm{tail}(e),e} = 1$,
$F_{\mathrm{head}(e),e} = -1$), the linear relaxation reads
\begin{subequations} \label{eq:crew}
\begin{align}
\min_{\xi_1,\dots,\xi_K}\quad & \sum_{k=1}^{K} c^{(k)\top} \xi_k\\
\text{s.t.}\quad & \Bigl(\sum_{k=1}^{K}A_{\mathrm{task}}\,\xi_k\Bigr)_i = 1 && (i\in\mathcal{R}_E),
\label{eq:crew-eq}\\
& \Bigl(\sum_{k=1}^{K}A_{\mathrm{task}}\,\xi_k\Bigr)_i \ge 1 && (i\in\mathcal{R}_I),
\label{eq:crew-ineq}\\
& F \xi_k = b^{(k)},\quad 0\le \xi_k\le 1 && (k = 1,\dots,K),
\end{align}
\end{subequations}
where $b^{(k)}$ has the entry $+1$ at the source, $-1$ at the sink and $0$ elsewhere, and
$\mathcal{R}_E$, $\mathcal{R}_I$ are the sets of equality and inequality coverage rows. This is a
multi-block linear program (\cref{def:multiblock}) with identical block polytopes. A
coverage row in $\mathcal{R}_E$ requires a flight to be operated exactly once; a row in
$\mathcal{R}_I$ requires it at least once, which allows a pilot to travel on a flight as a
passenger. The columns of the master problem are complete schedules, and the pricing
problem~\eqref{eq:pricing} of pilot $k$ is a shortest path problem with the edge weights
$c^{(k)} - A_{\mathrm{task}}^\top\pi$. These weights can be negative, but the graph is
acyclic, so one pass of dynamic programming in topological order solves it in
$O(|V|+|E|)$.

\subsection{Instance generator}

The instances are generated randomly but reproducibly.
\begin{itemize}
    \item \textbf{Network.} The nodes $1,\dots,|V|$ are ordered in time. A backbone edge
          connects every node to the next one, which guarantees a source--sink path. A skip
          edge from a node to any node at least two steps later exists independently with a
          fixed probability. An edge that spans $\ell$ steps costs $\ell + U(0,1)$, multiplied
          by $1 + 0.1\,U(0,1)$ for each pilot separately, where $U(0,1)$ is a uniform random
          number.
    \item \textbf{Flights.} Each skip edge operates a random set of up to ten flights whose
          position in time is close to that of the edge. A flight is only attached to edges
          that overlap in time, so that a single schedule operates every flight at most once.
          The task matrix therefore has a banded structure.
    \item \textbf{Feasibility.} Only flights that can be covered exactly once in the linear
          relaxation are kept, so that the equality variant is feasible.
    \item \textbf{Inequality rows.} In the inequality variant, a random 30\% of the coverage
          rows are turned into ``$\ge1$'' rows.
\end{itemize}

\Cref{tab:instances} lists the six instance sizes. The instances I1--I4 are used for the
comparison of all methods; S5 and S6, with up to 960 coverage rows, extend the study of
scaling. The last column is the number of variables of the full edge-based
formulation~\eqref{eq:crew}, which HiGHS solves directly to provide the reference optimum
$z^\star$.

\begin{table}[ht]
\centering
\renewcommand{\arraystretch}{1.2}
\begin{tabular}{lrrrrrr}
\toprule
instance & $|V|$ & $|E|$ & $K$ & $m$ & $m_I$ (inequality variant) & variables of~\eqref{eq:crew} \\
\midrule
\tabinput{data/tab_instances.tex}
\bottomrule
\end{tabular}
\caption{Test instances. $K$ is the number of pilots, $m$ the number of coverage rows, and
$m_I$ the number of ``$\ge$'' rows in the inequality variant.}
\label{tab:instances}
\end{table}


\section{Implementation and Setup}
\label{sec:setup}

All methods are implemented in Julia~1.12 \cite{bezanson2017julia} in one framework: the
column generation driver, the pricing algorithm, the column pool and the statistics are
shared, and only the solver of the restricted master problem differs. We compare six
methods:

\begin{center}
\renewcommand{\arraystretch}{1.2}
\resizebox{\textwidth}{!}{%
\begin{tabular}{ll}
\toprule
method & master problem \\
\midrule
Exact LP (HiGHS) & none: formulation~\eqref{eq:crew} solved directly by the HiGHS simplex method \\
LP CG (HiGHS) & classical linear master problem, HiGHS simplex, warm-started \\
Explicit QP CG (HiGHS) & explicit regularized master problem~\eqref{eq:explicit}, HiGHS QP solver \\
Eliminated QP CG (HiGHS) & eliminated regularized master problem~\eqref{eq:eliminated}, HiGHS QP \\
QPAS CG, sparse factor & eliminated regularized master problem, \cref{alg:hot-qpas} with the factorization of \cref{chap:updates} \\
QPAS CG, dense factor & eliminated regularized master problem, \cref{alg:hot-qpas} with the factorization of \cref{chap:dense} \\
\bottomrule
\end{tabular}
}
\end{center}

HiGHS \cite{huangfu2018parallelizing} (version 1.15, with its simplex method for the LPs and
its active-set QP solver for the QPs) is used through a persistent model, to which new columns
are added incrementally, so that its LP solves are warm-started. The LP master starts with
artificial columns of cost $M = 10^4$ on both sides of every linking row (\cref{sec:column-generation}),
which keep it feasible until enough columns are generated. The eliminated QP master is
run on the equality variant only. The sparse
factorization of the active-set method uses CHOLMOD \cite{chen2008cholmod} from SuiteSparse~7.10
for the factorization and its updates. The dense factorization is our own implementation of
\cref{chap:dense}; its triangular solves call the BLAS routine \code{trsv} of OpenBLAS, and
its rank-one modifications are plain loops over the columns of the stored array.

\medskip

All column generation runs use the same settings: the initial columns are the cheapest path
of every pilot, pricing visits the pilots in random order and adds the first improving
column, the proximal point rule moves the center after every iteration, and $\delta = 1$,
$\varepsilon_{\mathrm{c}} = 10^{-6}$, $\varepsilon_{\mathrm{rc}} = 10^{-6}$. The active-set
method uses the column capacity $\nmax = 20\,000$ and the pivot tolerance
$\varepsilon_{\mathrm{piv}} = 10^{-10}$ with both factorizations. The sparse factorization is
refactorized every $\tau = 150$ updates unless stated otherwise; the dense factorization is
only refactorized when the pivot invariant is violated (\cref{sec:dense-no-zeros}).

\medskip

The HiGHS QP masters are slow, so they are run to completion on I1 only; on the equality
variant of I2 they are stopped after 120 seconds, which is enough to compare their path and
their time per master problem with the active-set method. The LP master is run on I1--I3 with
a limit of 900 seconds, and the direct LP solve of S5 and S6 is stopped after 300 seconds. All
experiments ran single-threaded on an AMD Ryzen~7 7700 processor with 32\,GB of memory under
Windows~11. Times of the active-set method are the faster of two runs, except for runs longer
than one minute on S5 and S6, which were run once; all other times are from a single run.


\section{Correctness}
\label{sec:results-correctness}

\Cref{tab:comparison-eq,tab:comparison-ineq} compare all methods on both variants. Besides the
status, the number of column generation iterations, the number of generated columns and the
times, they report three accuracy measures relative to the reference optimum $z^\star$:
\begin{itemize}
    \item the relative error of the cost, $(c^\top\lambda - z^\star)/|z^\star|$;
    \item the relative gap of the best Lagrangian bound, $(z^\star - \theta)/|z^\star|$, which
          is nonnegative for every valid bound (\cref{prop:lagrangian-bound});
    \item the infeasibility $\lVert b - A\lambda\rVert_\infty$ of the final weights in the linking
          rows (for inequality rows only the violated part).
\end{itemize}

\begin{table}[!htbp]
\centering
\scriptsize
\renewcommand{\arraystretch}{1.15}
\setlength{\tabcolsep}{4pt}
\resizebox{\textwidth}{!}{%
\begin{tabular}{llcrrrrccc}
\toprule
 & method & status & CG iter. & columns & time [s] & master [ms/iter.] & cost error & bound gap & infeasibility \\
\midrule
\tabinput{data/tab_comparison_eq.tex}
\bottomrule
\end{tabular}
}
\caption{Equality variant: all methods. ``time'' denotes a method stopped by its time limit
(120\,s for the HiGHS QP masters on I2). Master times are averages per column generation
iteration.}
\label{tab:comparison-eq}
\end{table}

\begin{table}[!htbp]
\centering
\scriptsize
\renewcommand{\arraystretch}{1.15}
\setlength{\tabcolsep}{4pt}
\resizebox{\textwidth}{!}{%
\begin{tabular}{llcrrrrccc}
\toprule
 & method & status & CG iter. & columns & time [s] & master [ms/iter.] & cost error & bound gap & infeasibility \\
\midrule
\tabinput{data/tab_comparison_ineq.tex}
\bottomrule
\end{tabular}
}
\caption{Inequality variant (30\% ``$\ge$'' rows). The eliminated HiGHS QP master is not
run on this variant, and the explicit HiGHS QP master is run on I1 only.}
\label{tab:comparison-ineq}
\end{table}

\paragraph{All methods reach the optimum.} Every column generation run that finishes within
its time limit stops with an optimal status, and its Lagrangian bound agrees with $z^\star$
to within $2\cdot10^{-7}$ relative or better. This confirms the exact termination criterion
of \cref{thm:exact-termination} in practice, including the sign conditions of the inequality
rows: the bound is evaluated at the projected duals $\pi^+$, since at duals that violate them it
would not be valid and could exceed $z^\star$ (\cref{sec:bound-ineq,rem:bound-sign}).

\paragraph{The accuracy of the stabilized methods.} The stabilized methods stop when
$\lVert\pi-\pihat\rVert_\infty\le\varepsilon_{\mathrm{c}}\max(1,\lVert\pi\rVert_\infty)$, which by
\cref{prop:stop-accuracy} allows an infeasibility of up to $\delta\varepsilon_{\mathrm{c}}\max(1,\lVert\pi\rVert_\infty)$
in the linking rows. The infeasibility column shows exactly this order of magnitude, about
$10^{-5}$ for the larger instances, where the duals are of order $10$ to $100$. The cost
$c^\top\lambda$ of such a slightly infeasible solution can be off by more than the
infeasibility itself: on I2 it is $2\cdot10^{-5}$ below $z^\star$. The Lagrangian bound, which
certifies the optimal value, is unaffected and agrees with $z^\star$ to about $2\cdot10^{-7}$. A
tighter $\varepsilon_{\mathrm{c}}$ removes the cost error at the price of a few more center-only iterations. In a
test on a smaller instance with $\delta = 5$, the cost error dropped from $10^{-6}$ for
$\varepsilon_{\mathrm{c}} = 10^{-6}$ to $10^{-15}$ for $\varepsilon_{\mathrm{c}} = 10^{-9}$, with one additional iteration.

\paragraph{The active-set method and HiGHS solve the same problem.} On instance I1, where
all stabilized methods finish, the explicit and eliminated HiGHS QPs and the active-set method
perform almost the same number of column generation iterations and generate almost the same
number of columns: they follow the same sequence of master problems, as they should. The small
differences come from the accuracy of the master solves, which affects the duals passed to
pricing and hence the order in which columns are found.

\paragraph{The two factorizations solve the same problem.} The two variants of the active-set
method differ only in the linear algebra, so in exact arithmetic they produce the same iterates
(\cref{sec:dense-comparison}). In floating point they agree closely: their final objectives
differ by at most $6\cdot10^{-7}$ relative on all twelve runs (\cref{tab:backends} below),
with the same accuracy against $z^\star$, and on I2 they perform exactly the same 11\,974
iterations. Where the iteration counts differ, rounding differences have broken a tie in a
ratio test or a pricing round differently, after which the runs continue on slightly
different paths. The differences are below 7\% on all runs but one: on the inequality
variant of S6 the dense variant needs 9\,574 iterations against 6\,926, with the same final
objective.


\section[Convergence and the Two Phases]{Convergence and the Two Phases of Regularized Column Generation}
\label{sec:results-phases}

\Cref{fig:convergence} shows the convergence of the master objective on instance I2 for all
column generation methods, together with the change of the duals from one iteration to the
next. Here and in the rest of this section, ``QPAS'' denotes the active-set method with the
dense factorization; the sparse one follows the same path (it performs the same iterations on
I2).

\begin{figure}[ht]
\centering
\begin{tikzpicture}
\begin{groupplot}[group style={group size=2 by 1, horizontal sep=1.6cm},
    width=0.5\textwidth, height=0.42\textwidth, ymode=log,
    xlabel={CG iteration}, tick label style={font=\scriptsize}, label style={font=\small},
    title style={font=\small}, legend style={font=\scriptsize, at={(0.98,0.98)}, anchor=north east},
    every axis plot/.append style={line width=0.7pt}]
\nextgroupplot[title={relative error of the master objective}, ymin=1e-9, ymax=1]
    \addplot[cred] table[x=it, y=gap, col sep=comma] {data/conv_I2_0_lp.csv};
    \addplot[cgreen] table[x=it, y=gap, col sep=comma] {data/conv_I2_0_explicit_qp.csv};
    \addplot[cpurple, densely dashed] table[x=it, y=gap, col sep=comma] {data/conv_I2_0_eliminated_qp.csv};
    \addplot[cblue] table[x=it, y=gap, col sep=comma] {data/conv_I2_0_qpas_dense.csv};
    \legend{LP CG, explicit QP, eliminated QP, QPAS}
\nextgroupplot[title={dual change $\lVert\pi_t - \pi_{t-1}\rVert_\infty$}, ymin=1e-6]
    \addplot[cred, only marks, mark size=0.3pt] table[x=it, y=dpi, col sep=comma] {data/conv_I2_0_lp.csv};
    \addplot[cblue, only marks, mark size=0.3pt] table[x=it, y=dpi, col sep=comma] {data/conv_I2_0_qpas_dense.csv};
    \legend{LP CG, QPAS}
\end{groupplot}
\end{tikzpicture}
\caption{Instance I2 (equality variant). Left: relative distance of the master objective to
$z^\star$. Right: change of the duals between consecutive iterations. The two HiGHS QP masters
stop at their time limit of 120 seconds after about 2\,000 iterations.}
\label{fig:convergence}
\end{figure}

\paragraph{Stabilization.} The duals of the linear master problem jump by up to $2\cdot10^4$
in the first iterations, when the artificial columns of cost $M = 10^4$ are still active, and
by about $2$ in a typical iteration of the second half of the run. The regularized duals change
by about $10^{-4}$ in a typical iteration of the second half. The linear master problem,
however, reaches the optimum in fewer iterations on this instance. Its jumps explore the dual
space quickly, while the proximal term keeps the regularized duals close to the center, so
that they need many small steps. Regularization therefore does not reduce the number of
iterations here. What it buys is iterations whose master problems differ only slightly from
each other, which is exactly what the hot start exploits: an iteration of the active-set
method costs $0.08$\,ms on average with the dense factorization and $0.19$\,ms with the
sparse one, against $29$\,ms for the warm-started LP master (\cref{tab:comparison-eq}).

\paragraph{The QP masters follow the same path.} The curves of the explicit QP, the
eliminated QP and the active-set method are almost indistinguishable over the first 2\,000
iterations, after which the HiGHS QP masters run out of time.

\medskip

The right panel of \cref{fig:convergence} shows a second feature of the regularized methods:
from some point on, the duals change by very small amounts in most iterations. These are the
center-only iterations of \cref{sec:center-update}. \Cref{fig:phases} shows how the two kinds
of iterations are distributed over a run.

\begin{figure}[ht]
\centering
\begin{tikzpicture}
\begin{groupplot}[group style={group size=2 by 2, horizontal sep=1.6cm, vertical sep=0.6cm,
    xlabels at=edge bottom, xticklabels at=edge bottom},
    width=0.5\textwidth, height=0.3\textwidth, xlabel={CG iteration},
    tick label style={font=\scriptsize}, label style={font=\small}, title style={font=\small},
    every axis plot/.append style={line width=0.7pt}]
\nextgroupplot[title={I2, equality variant}, ylabel={columns}]
    \addplot[cblue] table[x=it, y=ncol, col sep=comma] {data/conv_I2_0_qpas_dense.csv};
\nextgroupplot[title={I2, inequality variant}]
    \addplot[cblue] table[x=it, y=ncol, col sep=comma] {data/conv_I2_30_qpas_dense.csv};
\nextgroupplot[ylabel={share adding a column}, ymin=0, ymax=1.05]
    \addplot[const plot mark mid, fill=cblue!35, draw=cblue] table[x=it, y=share, col sep=comma]
        {data/phases_I2_0.csv} \closedcycle;
\nextgroupplot[ymin=0, ymax=1.05]
    \addplot[const plot mark mid, fill=cblue!35, draw=cblue] table[x=it, y=share, col sep=comma]
        {data/phases_I2_30.csv} \closedcycle;
\end{groupplot}
\end{tikzpicture}
\caption{The two phases of regularized column generation (active-set method, instance I2).
Top: number of columns. Bottom: share of the iterations in each window of 100 iterations
that add a column; the remaining iterations are center-only iterations.}
\label{fig:phases}
\end{figure}

\paragraph{Two phases.} A run starts with a \textbf{column phase}, in which (almost) every
iteration adds a column. It is followed by a \textbf{proximal phase}, in which most
iterations are center-only iterations that move the center towards the restricted dual
optimum, interrupted by occasional iterations that find a new column because the duals have
moved. The phases are not strictly separated: columns keep appearing far into the proximal
phase.
\Cref{tab:phases} quantifies this for all instances.

\begin{table}[ht]
\centering
\small
\renewcommand{\arraystretch}{1.15}
\setlength{\tabcolsep}{4pt}
\resizebox{\textwidth}{!}{%
\begin{tabular}{lcrrrrrrrr}
\toprule
 & & & & & first & last & \multicolumn{2}{c}{mean inner iter.} & master time \\
instance & ineq. & CG iter. & with column & center-only & center-only & column & column & center-only & center-only \\
\midrule
\tabinput{data/tab_phases.tex}
\bottomrule
\end{tabular}
}
\caption{Column-adding and center-only iterations of the active-set method (dense
factorization). ``inner iter.'' counts the iterations of the active-set method per master
problem.}
\label{tab:phases}
\end{table}

The two kinds of iterations differ sharply in cost. A column-adding iteration needs on
average tens of active-set iterations, since the new column is released and the working set
adapts to it. A center-only iteration needs at most two on average: as shown in
\cref{prop:center-only}, once the working set has settled, the direction is zero, and the
master problem is solved by a single multiplier test. The length of the proximal phase
varies a lot between instances. On I2 and I3 (equality variant) center-only iterations are
the large majority, on I4 they are a minority. On all instances except the tiny I1 their share
of the master time is much smaller than their share of the iterations.

\Cref{fig:optimality} shows how the optimality measures evolve over the two phases for the
active-set method on I2. In the column phase, the cost error, the error of the master
objective and the infeasibility decrease together. In the proximal phase, the Lagrangian
bound becomes available, since all pilots are priced in every center-only iteration. The
bound approaches $z^\star$ from below at the same rate at which the master objective
approaches it, so the two errors are indistinguishable in the plot. The infeasibility, and
with it the cost error, decrease more slowly and in steps: each step is a change of the
working set that ends one of the dual walks of \cref{prop:center-only}.

\begin{figure}[ht]
\centering
\begin{tikzpicture}
\begin{axis}[width=0.78\textwidth, height=0.42\textwidth, ymode=log, xlabel={CG iteration},
    ymin=1e-9, ymax=10, tick label style={font=\scriptsize}, label style={font=\small},
    legend style={font=\scriptsize, at={(0.98,0.98)}, anchor=north east, legend columns=2},
    every axis plot/.append style={line width=0.7pt}]
    \addplot[cred!45, only marks, mark=*, mark size=0.7pt] table[x=it, y=bgap, col sep=comma] {data/bound_I2_0_qpas_dense.csv};
    \addplot[corange] table[x=it, y=costgap, col sep=comma] {data/conv_I2_0_qpas_dense.csv};
    \addplot[cgreen] table[x=it, y=infeas, col sep=comma] {data/conv_I2_0_qpas_dense.csv};
    \addplot[cblue, densely dashed] table[x=it, y=gap, col sep=comma] {data/conv_I2_0_qpas_dense.csv};
    \legend{bound gap $(z^\star-\theta)/|z^\star|$, cost error $|c^\top\lambda - z^\star|/|z^\star|$,
            infeasibility, master objective error}
\end{axis}
\end{tikzpicture}
\caption{Optimality measures of the active-set method over a run (I2, equality variant). The
Lagrangian bound is only available in iterations that price all pilots, which are the
iterations that find no column.}
\label{fig:optimality}
\end{figure}


\section{Computational Time}
\label{sec:results-time}

\paragraph{Master time per iteration.} \Cref{fig:master-time} shows the time of the master
solve per column generation iteration on I2. The time of the LP master and of the two HiGHS QP
masters grows with the number of columns. The HiGHS QP solvers take 80--95\,ms per iteration
in the second half of their run, since they rebuild their internal factorization for every
master problem. The active-set method stays between $0.003$ and $1.6$\,ms per iteration with
the dense factorization and between $0.006$ and $6.5$\,ms with the sparse one, with medians
of $0.04$ and $0.07$\,ms in the second half of the run. That is about three orders of
magnitude below the HiGHS QP masters and almost three orders of magnitude below the median of
$30$\,ms of the warm-started LP master.

\begin{figure}[ht]
\centering
\begin{tikzpicture}
\begin{axis}[width=0.78\textwidth, height=0.42\textwidth, ymode=log, xlabel={CG iteration},
    ylabel={master time per iteration [ms]}, tick label style={font=\scriptsize},
    label style={font=\small}, legend style={font=\scriptsize, at={(0.98,0.5)}, anchor=east},
    every axis plot/.append style={line width=0.8pt}]
    \addplot[cred, dashdotted] table[x=it, y=mtavg, col sep=comma] {data/conv_I2_0_lp.csv};
    \addplot[cgreen] table[x=it, y=mtavg, col sep=comma] {data/conv_I2_0_explicit_qp.csv};
    \addplot[cpurple, densely dashed] table[x=it, y=mtavg, col sep=comma] {data/conv_I2_0_eliminated_qp.csv};
    \addplot[corange] table[x=it, y=mtavg, col sep=comma] {data/conv_I2_0_qpas.csv};
    \addplot[cblue] table[x=it, y=mtavg, col sep=comma] {data/conv_I2_0_qpas_dense.csv};
    \legend{LP CG, explicit QP, eliminated QP, QPAS (sparse), QPAS (dense)}
\end{axis}
\end{tikzpicture}
\caption{Master time per column generation iteration on I2 (moving average over 50
iterations).}
\label{fig:master-time}
\end{figure}

\paragraph{Where the time goes.} \Cref{fig:breakdown} splits the time of the active-set
method into its operations, for both factorizations, as shares of the total time of the run
with the sparse factorization. With the sparse factorization, the largest share, 32--48\%, is
the direction computation (the right-hand side and the sparse solve with the factor), followed
by the updates of the factorization and the periodic refactorization, together 26--55\%; the
refactorization share grows with the instance size. The dense factorization cuts the
direction computation by a factor of 2.5 to 4, the updates by a factor of 3 to 4.5, and the
refactorizations almost entirely. What remains is dominated by operations that do not touch
the factor: the multiplier test, the only operation whose cost grows with the number of
columns, is now about as expensive as the direction computation on the equality variants and
takes between a fifth and two fifths of the time of the dense runs. Pricing takes less than
7\% of the total time on all instances shown.

\begin{figure}[ht]
\centering
\begin{tikzpicture}
\begin{axis}[xbar stacked, table/col sep=comma, width=0.8\textwidth, height=0.95\textwidth, xmin=0, xmax=100,
    xlabel={share of the total time of the sparse run [\%]}, ytick=data,
    yticklabels from table={data/breakdown.csv}{label},
    y dir=reverse, bar width=6.5pt, tick label style={font=\scriptsize}, label style={font=\small},
    legend style={font=\scriptsize, at={(0.5,-0.08)}, anchor=north, legend columns=3},
    reverse legend=false]
    \addplot[fill=cblue!70, draw=cblue] table[x=direction, y expr=\coordindex, col sep=comma] {data/breakdown.csv};
    \addplot[fill=corange!70, draw=corange] table[x=multipliers, y expr=\coordindex, col sep=comma] {data/breakdown.csv};
    \addplot[fill=cgrey!50, draw=cgrey] table[x=step, y expr=\coordindex, col sep=comma] {data/breakdown.csv};
    \addplot[fill=cgreen!60, draw=cgreen] table[x=updates, y expr=\coordindex, col sep=comma] {data/breakdown.csv};
    \addplot[fill=cred!60, draw=cred] table[x=refactor, y expr=\coordindex, col sep=comma] {data/breakdown.csv};
    \addplot[fill=cpurple!50, draw=cpurple] table[x=pricing, y expr=\coordindex, col sep=comma] {data/breakdown.csv};
    \legend{direction, multipliers, step, updates, refactorization, pricing}
\end{axis}
\end{tikzpicture}
\caption{Time breakdown of the active-set method and of pricing with the sparse and the dense
factorization, as shares of the total time of the sparse run on the same instance. ``updates''
are the expansions and reductions of the working set without the refactorizations.}
\label{fig:breakdown}
\end{figure}

\Cref{tab:counters} lists the events inside the active-set method for every run. Over a full
run, the working set changes tens to hundreds of thousands of times, and every change is a
single update of the factorization. Representative swaps are rare, because representatives
carry the largest weights. Singular reduced Hessians occur, but only a handful of times per
run, and each is resolved by one zero-curvature step. The final dimension $|\Nx|$ of $\Heff$
stays below the number of coverage rows, as guaranteed by \cref{cor:dimension}. The sparse
factorization is rebuilt hundreds to thousands of times per run, almost all of them periodic
refactorizations; the dense one is rebuilt at most 16 times, each time because a singular
release had to be undone or a pivot had decayed.

\begin{table}[ht]
\centering
\scriptsize
\renewcommand{\arraystretch}{1.15}
\setlength{\tabcolsep}{3.5pt}
\resizebox{\textwidth}{!}{%
\begin{tabular}{lcrrrrrcrrrrr}
\toprule
 & & AS iter. & steps (full) & expans. & swaps & reduct. & surplus exp./red. & sing. & updates & refact. & $|\Nx|$ & $|\Ns|$ \\
\midrule
\tabinput{data/tab_qpas_counters.tex}
\bottomrule
\end{tabular}
}
\caption{Events inside the active-set method (dense factorization) over a whole column
generation run: active-set iterations, primal steps (of which full steps), expansions,
representative swaps, reductions, expansions/reductions of surpluses, singular steps,
factorization updates, refactorizations (sparse / dense factorization), and the final sizes of
$\Nx$ and $\Ns$.}
\label{tab:counters}
\end{table}

\paragraph{Scaling.} \Cref{fig:scaling} shows the total time of the active-set method with
both factorizations and of the direct LP solve as the instances grow, up to S6 with 960
coverage rows, and \cref{tab:scaling} gives the details of the two largest instances. From I2
to S6 the time of the dense variant grows almost exactly with the third power of $m$. Of this,
the number of active-set iterations accounts for a factor of about $m^2$ (from 62\,000 on I2 to
950\,000 on S6, \cref{tab:counters}) and the time per iteration for a factor of about $m$. The
latter is less than the $O(|\Nx|^2)$ of the factor operations because, at these sizes, the
parts of an iteration that grow only linearly, such as the right-hand side, the step and the
multiplier test, still take a large share of its time. The sparse variant grows faster, about
with $m^{3.5}$, because its refactorizations take an increasing share of the time, and the
direct LP solve grows faster still. With the dense
factorization, the complete column generation run is faster than the direct solution of the
full LP from I2 on in the equality variant and on all instances in the inequality variant, by
a factor of three to four on I4 and S5. On S6, and on the inequality variant of S5, the direct
LP solve did not finish within 300 seconds, while the active-set method with the dense
factorization solved both variants of S6 in 84--95 seconds. The inequality variant is
consistently faster than the equality variant, since it needs fewer active-set iterations
(\cref{sec:results-ineq}).

\begin{figure}[ht]
\centering
\begin{tikzpicture}
\begin{axis}[width=0.72\textwidth, height=0.5\textwidth, xmode=log, ymode=log,
    xlabel={coverage rows $m$}, ylabel={total time [s]}, tick label style={font=\scriptsize},
    label style={font=\small}, legend style={font=\scriptsize, at={(0.02,0.98)}, anchor=north west},
    every axis plot/.append style={line width=0.8pt}, log basis x=10, unbounded coords=jump]
    \addplot[corange, mark=triangle*, mark size=1.8pt] table[x=m, y=qpas, col sep=comma] {data/scaling_eq.csv};
    \addplot[corange, dashed, mark=triangle, mark size=1.8pt] table[x=m, y=qpas, col sep=comma] {data/scaling_ineq.csv};
    \addplot[cblue, mark=*, mark size=1.5pt] table[x=m, y=dense, col sep=comma] {data/scaling_eq.csv};
    \addplot[cblue, dashed, mark=o, mark size=1.5pt] table[x=m, y=dense, col sep=comma] {data/scaling_ineq.csv};
    \addplot[black!75, mark=square*, mark size=1.5pt] table[x=m, y=exact, col sep=comma] {data/scaling_eq.csv};
    \addplot[black!75, dashed, mark=square, mark size=1.5pt] table[x=m, y=exact, col sep=comma] {data/scaling_ineq.csv};
    \addplot[cgrey, dotted, domain=40:1000, samples=2] {1.1e-7*x^3};
    \legend{QPAS sparse (equality), QPAS sparse (inequality), QPAS dense (equality),
            QPAS dense (inequality), exact LP (equality), exact LP (inequality), $\propto m^3$}
\end{axis}
\end{tikzpicture}
\caption{Total time against the number of coverage rows for instances I1--I4, S5 and S6. The
exact LP of S6 and of the inequality variant of S5 did not finish within 300 seconds.}
\label{fig:scaling}
\end{figure}

\begin{table}[ht]
\centering
\scriptsize
\renewcommand{\arraystretch}{1.15}
\setlength{\tabcolsep}{3.5pt}
\resizebox{\textwidth}{!}{%
\begin{tabular}{lccccrrrrrrcc}
\toprule
 & ineq. & exact LP [s] & factor & status & CG iter. & columns & time [s] & master [ms/iter.] & refact. & $|\Nx|$ & bound gap & infeasibility \\
\midrule
\tabinput{data/tab_scaling.tex}
\bottomrule
\end{tabular}
}
\caption{The large instances S5 and S6. The bound gap is only available where the direct LP
solve finished within its limit of 300 seconds.}
\label{tab:scaling}
\end{table}


\section{Inequality Rows}
\label{sec:results-ineq}

With 30\% inequality rows, column generation needs fewer active-set iterations on all
instances, and fewer column generation iterations on I1--I4, considerably fewer on I1--I3
(\cref{tab:comparison-eq,tab:comparison-ineq,tab:counters}): over-covering a flight is allowed, so the
optimal schedules are found with fewer columns, and the proximal phase is shorter.

\medskip

\Cref{fig:sets} follows the sizes of the nullspace sets over a run with inequality rows.
Within the first iterations, about 35 surpluses are released: the initial schedules
over-cover their rows, so the duals of these rows become negative. These rows drop out of
$\Aeff$, and by \cref{cor:dimension} they also reduce the room for nullspace columns,
$|\Nx|\le m - |\Ns|$. Over the rest of the run $|\Ns|$ only changes slowly, between 37 and 51,
ending at 48, while individual surpluses are released and fixed again thousands of times
(\cref{tab:counters}). The number of nullspace columns grows during the column phase and ends
at $|\Nx| = 306$, just below its bound $m - |\Ns| = 308$: at the end of the run, the effective
reduced Hessian is almost as large as it can be.

\begin{figure}[ht]
\centering
\begin{tikzpicture}
\begin{axis}[width=0.78\textwidth, height=0.42\textwidth, xlabel={CG iteration},
    ylabel={set size}, tick label style={font=\scriptsize}, label style={font=\small},
    legend style={font=\scriptsize, at={(0.98,0.5)}, anchor=east},
    every axis plot/.append style={line width=0.8pt}, ymin=0]
    \addplot[cgrey, dashed] table[x=it, y=m, col sep=comma] {data/sets_I3_30.csv};
    \addplot[cred] table[x=it, y=bound, col sep=comma] {data/sets_I3_30.csv};
    \addplot[cblue] table[x=it, y=N, col sep=comma] {data/sets_I3_30.csv};
    \addplot[corange] table[x=it, y=Ns, col sep=comma] {data/sets_I3_30.csv};
    \legend{$m$, $m - |\Ns|$ (bound on $|\Nx|$), $|\Nx|$, $|\Ns|$}
\end{axis}
\end{tikzpicture}
\caption{Nullspace columns and free surpluses over a run (I3, inequality variant, dense
factorization), sampled every 20 iterations.}
\label{fig:sets}
\end{figure}

\medskip

Surplus operations are frequent (\cref{tab:counters}): with inequality rows, 14--19\% of all
expansions and reductions of the working set on I2--S6 are expansions and reductions of
surpluses, each of which is a rank-one update or downdate of the factorization. The average
time per master problem of the inequality variant is of the same order as that of the equality
variant (\cref{tab:comparison-eq,tab:comparison-ineq}). The remaining differences come from the
mix of iterations. On I2, for instance, the inequality variant has relatively few of the cheap
center-only iterations (\cref{tab:phases}), which raises its average. Singular reduced
Hessians, of columns or surpluses, occurred at most six times per run on I1--I4. The singular
surplus release of \cref{thm:zc-surplus} is exercised by the test suite of the implementation
on constructed instances.


\section{Fill-In and Refactorization of the Sparse Factor}
\label{sec:results-fill}

The factorization of $\Heff$ is updated tens of thousands of times per run
(\cref{tab:counters}). \Cref{fig:fill} shows what this does to the contents of the sparse
factor, using the quantities of \cref{def:stored-effective}, without refactorization, with
refactorization every 150 updates, and, for comparison, for the dense factor of
\cref{chap:dense}.

\begin{figure}[ht]
\centering
\begin{tikzpicture}
\begin{groupplot}[group style={group size=3 by 1, horizontal sep=0.5cm, yticklabels at=edge left},
    width=0.38\textwidth, height=0.4\textwidth, xlabel={CG iteration}, ymode=log,
    ymin=100, ymax=2e6, tick label style={font=\scriptsize}, label style={font=\small},
    title style={font=\small}, legend style={font=\scriptsize, at={(0.98,0.02)}, anchor=south east},
    every axis plot/.append style={line width=0.7pt}]
\nextgroupplot[title={sparse, never refactorized}, ylabel={off-diagonal entries of $L$}]
    \addplot[cred] table[x=it, y=stored, col sep=comma] {data/fill_I3_tau0.csv};
    \addplot[cblue] table[x=it, y=effective, col sep=comma] {data/fill_I3_tau0.csv};
    \addplot[cgrey, dashed] table[x=it, y=dense, col sep=comma] {data/fill_I3_tau0.csv};
    \legend{stored, nonzero, dense equiv.}
\nextgroupplot[title={sparse, $\tau = 150$}]
    \addplot[cred] table[x=it, y=stored, col sep=comma] {data/fill_I3_tau150.csv};
    \addplot[cblue] table[x=it, y=effective, col sep=comma] {data/fill_I3_tau150.csv};
    \addplot[cgrey, dashed] table[x=it, y=dense, col sep=comma] {data/fill_I3_tau150.csv};
\nextgroupplot[title={dense factor}]
    \addplot[cred] table[x=it, y=stored, col sep=comma] {data/fill_I3_dense_tau0.csv};
    \addplot[cblue] table[x=it, y=effective, col sep=comma] {data/fill_I3_dense_tau0.csv};
    \addplot[cgrey, dashed] table[x=it, y=dense, col sep=comma] {data/fill_I3_dense_tau0.csv};
\end{groupplot}
\end{tikzpicture}
\caption{Contents of the factor $L$ of $\Heff$ over a run (I3, equality variant), sampled
every 25 iterations: stored off-diagonal entries, the nonzero ones among them, and the
off-diagonal entries of a dense factor of the same dimension. In the right panel the three
curves coincide.}
\label{fig:fill}
\end{figure}

\paragraph{The true factor is dense.} In all panels the nonzero entries coincide with the
dense equivalent: the factor of $\Heff$ has essentially no exploitable sparsity. This confirms
\cref{rem:natural-ordering} and the starting point of \cref{chap:dense}: a fill-reducing
ordering cannot help, and the cost of a solve grows quadratically with $|\Nx|$.

\paragraph{Structural zeros accumulate.} Without the periodic refactorization, the stored
pattern of the sparse factor only grows. On I3 it reaches almost a million entries in the
column phase, about thirty times the size of the actual factor, and almost all of them are
structural zeros. Every solve and every update pays for them, and they connect active indices
with inactive ones, which enlarges the reach of the sparse solves. The drop around iteration
1\,400 is a refactorization triggered by a degraded pivot (\cref{sec:detection}). It resets the
pattern, after which it grows again and stays at 1.2 to 2.6 times the size of the factor. With
refactorization every 150 updates, the stored pattern stays within a small margin of the
effective one for most of the run. The dense factor stores exactly the dense equivalent at
every moment (\cref{prop:no-zeros}).

\paragraph{The refactorization interval.} \Cref{tab:refactor} and \cref{fig:refactor} vary the
refactorization interval of the sparse factorization. Without refactorization, the method is
three to five times slower, and the sparse factor needs six to eleven times more memory.
Refactorizing too often is also expensive, since every rebuild costs $O(|\Nx|^3)$
(\cref{prop:refactor-cost}): with an interval of 25, I3 needs almost 3\,000 rebuilds and is
70\% slower than with the default. Between these extremes the optimum is flat: any interval
between 100 and 600 is within about 10\% of the best time, and the default of 150 is within
3\% of the best on both instances. For the dense factorization, which has no structural zeros,
a periodic refactorization only costs time: with $\tau = 150$ it is 5--10\% slower than without.
Both dense variants are faster than the best sparse one by a factor of 2.1 to 2.6.

\begin{table}[ht]
\centering
\small
\renewcommand{\arraystretch}{1.15}
\resizebox{\textwidth}{!}{%
\begin{tabular}{lllrrrrrrr}
\toprule
 & factor & interval $\tau$ & time [s] & refact. & max.\ stored & max.\ zeros & mean zero share & reach$/|\Nx|$ & peak memory [MB] \\
\midrule
\tabinput{data/tab_refactor.tex}
\bottomrule
\end{tabular}
}
\caption{Effect of the refactorization interval (equality variants). ``max.\ zeros'' is the
largest number of stored entries that are zero; for the dense factor these are entries of
$L$ that happen to vanish, not structural zeros. ``reach$/|\Nx|$'' is the average number of
rows touched by a solve per nullspace column. ``never'' still refactorizes when a pivot
degrades. These runs form a separate series (times are the faster of two runs), so the times
differ slightly from \cref{tab:comparison-eq}.}
\label{tab:refactor}
\end{table}

\begin{figure}[ht]
\centering
\begin{tikzpicture}
\begin{axis}[width=0.68\textwidth, height=0.42\textwidth, xmode=log, xlabel={refactorization interval $\tau$},
    ylabel={total time [s]}, xtick={25,50,100,150,300,600,2000},
    xticklabels={25,50,100,150,300,600,never}, tick label style={font=\scriptsize},
    label style={font=\small}, legend style={font=\scriptsize, at={(0.02,0.98)}, anchor=north west},
    ymin=0, unbounded coords=jump, every axis plot/.append style={line width=0.8pt}]
    \addplot[cblue, mark=*, mark size=1.5pt] table[x=tau, y=time, col sep=comma] {data/refactor_I2.csv};
    \addplot[corange, mark=square*, mark size=1.5pt] table[x=tau, y=time, col sep=comma] {data/refactor_I3.csv};
    \addplot[cblue, dashed, domain=20:2500] table[x=tau, y=time, col sep=comma] {data/refactor_I2_dense.csv};
    \addplot[corange, dashed] table[x=tau, y=time, col sep=comma] {data/refactor_I3_dense.csv};
    \legend{I2 sparse, I3 sparse, I2 dense, I3 dense}
\end{axis}
\end{tikzpicture}
\caption{Total time of the active-set method against the refactorization interval of the
sparse factorization. The dashed lines show the dense factorization without periodic
refactorization.}
\label{fig:refactor}
\end{figure}


\section{The Sparse and the Dense Factorization}
\label{sec:results-factorizations}

\Cref{tab:backends} compares the two factorizations on all instances and variants. Both
variants of the active-set method reach the same objective and the same accuracy
(\cref{sec:results-correctness}), and their iteration counts are close, so the time
difference is caused by the linear algebra.

\begin{table}[ht]
\centering
\scriptsize
\renewcommand{\arraystretch}{1.15}
\setlength{\tabcolsep}{4pt}
\resizebox{\textwidth}{!}{%
\begin{tabular}{lcrrrrrrrrrrc}
\toprule
 & & & CG iter. & \multicolumn{3}{c}{master time [s]} & \multicolumn{2}{c}{total time [s]} & refact. & \multicolumn{2}{c}{factor memory [MB]} & objective \\
\cmidrule(lr){5-7}\cmidrule(lr){8-9}\cmidrule(lr){11-12}
instance & ineq. & $m$ & sparse / dense & sparse & dense & speedup & sparse & dense & sparse / dense & sparse & dense & rel.\ difference \\
\midrule
\tabinput{data/tab_backends.tex}
\bottomrule
\end{tabular}
}
\caption{The two factorizations of the effective reduced Hessian. Master times are summed over
the run, the speedup is their ratio, and the factor memory is the peak of the memory held by
the factorization.}
\label{tab:backends}
\end{table}

\paragraph{Speed.} The dense factorization makes the master problems 2.3 to 4.2 times
faster, and the speedup grows with the number of linking rows: 2.4 to 2.8 at $m = 225$ and $356$,
3.3 at $m = 480$, 3.6 at $m = 668$ and 4.2 at $m = 960$ (equality variants). Two effects
contribute. First, every operation on the factor is cheaper, as shown below. Second, the sparse
factorization has to be rebuilt periodically, and since a rebuild costs $O(|\Nx|^3)$ while the
interval stays fixed, these rebuilds take a growing share of its time
(\cref{fig:breakdown}): almost 40\% on S6, with more than 4\,000 rebuilds. The dense
factorization was rebuilt at most 16 times per run.

\paragraph{Cost per operation.} \Cref{tab:ops} lists the average time of the individual
factor operations as measured inside the runs. On the instances of the study, a solve and a
row addition are about five times faster with the dense factorization, and a row deletion five
to seven times. Assembling
a new row of $\Heff$ from the column pool, which is the same code for both factorizations,
now costs more than adding it to the factor.

\begin{table}[ht]
\centering
\small
\renewcommand{\arraystretch}{1.15}
\begin{tabular}{lrrrrrrrr}
\toprule
 & \multicolumn{2}{c}{solve} & \multicolumn{2}{c}{row addition} & \multicolumn{2}{c}{row deletion} & \multicolumn{2}{c}{assembly of a row} \\
\cmidrule(lr){2-3}\cmidrule(lr){4-5}\cmidrule(lr){6-7}\cmidrule(lr){8-9}
instance & sparse & dense & sparse & dense & sparse & dense & sparse & dense \\
\midrule
\tabinput{data/tab_ops.tex}
\bottomrule
\end{tabular}
\caption{Average time per factor operation in $\mu$s over a column generation run (equality
variants). Row additions include those of the refactorizations.}
\label{tab:ops}
\end{table}

To separate the cost of the operations from the behaviour of the active-set method,
\cref{fig:micro} and \cref{tab:micro} measure them in isolation. The matrix is the Gram matrix
of $k+1$ random sparse columns with entries in $\{-1,0,1\}$, which is dense like $\Heff$, both
factorizations hold its first $k$ indices, and every operation is repeated and the median time
is reported.

\begin{figure}[ht]
\centering
\begin{tikzpicture}
\begin{groupplot}[group style={group size=2 by 1, horizontal sep=1.4cm},
    width=0.5\textwidth, height=0.42\textwidth, xmode=log, ymode=log, xlabel={active size $k$},
    tick label style={font=\scriptsize}, label style={font=\small}, title style={font=\small},
    legend style={font=\scriptsize, at={(0.02,0.98)}, anchor=north west},
    every axis plot/.append style={line width=0.8pt}]
\nextgroupplot[title={solve and row addition}, ylabel={time per operation [$\mu$s]}]
    \addplot[corange, mark=*, mark size=1.2pt] table[x=k, y=sparse, col sep=comma] {data/micro_solve.csv};
    \addplot[cblue, mark=*, mark size=1.2pt] table[x=k, y=dense, col sep=comma] {data/micro_solve.csv};
    \addplot[corange, dashed, mark=square*, mark size=1.2pt] table[x=k, y=sparse, col sep=comma] {data/micro_append.csv};
    \addplot[cblue, dashed, mark=square*, mark size=1.2pt] table[x=k, y=dense, col sep=comma] {data/micro_append.csv};
    \legend{solve (sparse), solve (dense), append (sparse), append (dense)}
\nextgroupplot[title={deletion and rank-one update}]
    \addplot[corange, mark=*, mark size=1.2pt] table[x=k, y=sparse, col sep=comma] {data/micro_delete_first.csv};
    \addplot[cblue, mark=*, mark size=1.2pt] table[x=k, y=dense, col sep=comma] {data/micro_delete_first.csv};
    \addplot[corange, dashed, mark=square*, mark size=1.2pt] table[x=k, y=sparse, col sep=comma] {data/micro_update.csv};
    \addplot[cblue, dashed, mark=square*, mark size=1.2pt] table[x=k, y=dense, col sep=comma] {data/micro_update.csv};
    \legend{delete oldest (sparse), delete oldest (dense), update (sparse), update (dense)}
\end{groupplot}
\end{tikzpicture}
\caption{Median time of the factor operations against the number $k$ of active indices, for a
dense Gram matrix. The sparse factor has the dimension $\nmax = 20\,000$.}
\label{fig:micro}
\end{figure}

\begin{table}[ht]
\centering
\small
\renewcommand{\arraystretch}{1.15}
\resizebox{\textwidth}{!}{%
\begin{tabular}{lrrrrrrrr}
\toprule
 & \multicolumn{2}{c}{$k = 100$} & \multicolumn{2}{c}{$k = 350$} & \multicolumn{2}{c}{$k = 1000$} & \multicolumn{2}{c}{$k = 2000$} \\
\cmidrule(lr){2-3}\cmidrule(lr){4-5}\cmidrule(lr){6-7}\cmidrule(lr){8-9}
operation & sparse & dense & sparse & dense & sparse & dense & sparse & dense \\
\midrule
\tabinput{data/tab_micro.tex}
\bottomrule
\end{tabular}
}
\caption{Median time per operation in $\mu$s, and the memory of the factor, against the number
$k$ of active indices.}
\label{tab:micro}
\end{table}

The dense factorization is faster for every operation and every size. The largest gains are
for the operations that the sparse factorization performs through index structures: solves
and appends are 2 to 10 times faster, refactorizations 6 to 12 times. Deletions and rank-one
modifications, which both factorizations perform by the same column-oriented rank-one
recursion, are 1.2 to 2.6 times faster. For the sparse factorization the time per operation
grows faster than $k^2$ beyond $k\approx1000$, the dense one stays close to $k^2$. The memory
of the dense factor is seven to nine times smaller, since it neither stores indices nor embeds
the factor in $\nmax$ dimensions.

\paragraph{Deletions.} By \cref{prop:delete-cost}, the cost of a deletion depends on the
position of the removed index. Over a run, the trailing part $(t-\ell)/t$ that a deletion
updates was on average 0.22 to 0.36 of the factor on all instances: as anticipated in
\cref{sec:dense-delete}, most removed indices had been active only for a short time. The
dense factor never needed more than $m+1$ positions, as guaranteed by~\eqref{eq:dense-size} and
the one index of an undone singular release (\cref{sec:why-dense}).

\paragraph{Accuracy.} \Cref{fig:residual} shows the residual
$\max|(LDL^\top)_{\Nx\Nx} - \Heff|/\max(1,\max|\Heff|)$ of the factor over a run, with $\Heff$
recomputed from the columns at every sample. For both factorizations it stays at the level of
the rounding error, below $3\cdot10^{-15}$ on all runs of I1--I4, and it does not grow over
the run. The dense factor is rebuilt only a handful of times per run and yet does not drift,
even over almost 150\,000 updates on I4: the rank-one modifications of \cref{alg:c1} are
stable in practice, as \cref{rem:c1-stability} suggests.

\begin{figure}[ht]
\centering
\begin{tikzpicture}
\begin{groupplot}[group style={group size=2 by 1, horizontal sep=1.4cm, yticklabels at=edge left},
    width=0.5\textwidth, height=0.38\textwidth, ymode=log, ymin=1e-17, ymax=1e-13,
    xlabel={CG iteration}, tick label style={font=\scriptsize}, label style={font=\small},
    title style={font=\small}, legend style={font=\scriptsize, at={(0.98,0.98)}, anchor=north east},
    every axis plot/.append style={line width=0.7pt}]
\nextgroupplot[title={I3, equality variant}, ylabel={relative residual}]
    \addplot[corange] table[x=it, y=residual, col sep=comma] {data/resid_I3_0_sparse.csv};
    \addplot[cblue] table[x=it, y=residual, col sep=comma] {data/resid_I3_0_dense.csv};
    \legend{sparse, dense}
\nextgroupplot[title={I4, inequality variant}]
    \addplot[corange] table[x=it, y=residual, col sep=comma] {data/resid_I4_30_sparse.csv};
    \addplot[cblue] table[x=it, y=residual, col sep=comma] {data/resid_I4_30_dense.csv};
\end{groupplot}
\end{tikzpicture}
\caption{Relative residual of the factorization over a run, sampled every 20 iterations.}
\label{fig:residual}
\end{figure}


\section{Discussion}
\label{sec:discussion}

The experiments support the main claims of this work.
\begin{itemize}
    \item \textbf{Correctness.} The active-set method solves the same sequence of master
          problems as a general-purpose QP solver, reproduces the optimum of the full linear
          program, and certifies it by the Lagrangian bound, for equality and inequality rows
          and with both factorizations.
    \item \textbf{Speed.} With the dense factorization, a master problem of the active-set
          method costs between $0.003$ and $13$\,ms on average on our instances, and $0.08$\,ms
          on I2. That is almost three orders of magnitude less than a HiGHS QP master and two
          to three orders of magnitude less than a warm-started LP master. This makes regularized column
          generation much faster than classical column generation, although it often needs
          more iterations: on I3, 3.0\,s against 663\,s. The complete run is faster than the
          direct solution of the full LP from I2 on, and on the largest instances the direct
          solve does not finish within the time in which the active-set method is done.
    \item \textbf{Inequality rows} add no new kind of linear algebra: the surplus operations
          are rank-one modifications of the same factorization, free surpluses lower the bound
          $m - |\Ns|$ on the size of $\Heff$, the time per master problem stays of the same
          order, and the runs are shorter.
    \item \textbf{The factorization.} The factor of $\Heff$ is dense. A sparse factorization
          needs periodic refactorization to control its structural zeros, and a fixed interval
          of a few hundred updates works well. A dense factorization in activation order
          needs none, is two to four times faster with a speedup that grows with $m$, needs an
          order of magnitude less memory, and is just as accurate.
\end{itemize}

The experiments also show the limits of the approach. Since $\Heff$ is dense, the cost of the
method grows with the cube of the number of linking rows, and the \textbf{proximal phase} can
be long, since the duals approach the restricted optimum in many small steps
(\cref{prop:center-only}). With the dense factorization the factor is no longer the dominant
cost: the multiplier test, which touches every column of the working set, now costs as much as
the direction computation. Natural next steps are therefore an incremental multiplier test that
avoids rescanning all columns, and an adaptive stabilization parameter in the proximal phase.
A change of $\delta$ does not affect the factorization (\cref{rem:factor-independent}) and is
therefore free for this method.
